%!LW recipe=pdflatex x3
\documentclass[10pt,a4paper]{article}

\usepackage[spanish,es-nodecimaldot]{babel}
\usepackage[utf8]{inputenc}
\usepackage[T1]{fontenc}
\usepackage{newtxtext,newtxmath}
\usepackage{geometry}
\usepackage{microtype}
\usepackage{amsmath,mathtools}
\usepackage{enumitem}
\usepackage{xcolor}
\usepackage{hyperref}
\usepackage{fancyhdr}
\usepackage{lastpage}
\usepackage[most]{tcolorbox}

\geometry{margin=2.5cm,includeheadfoot}
\setlength{\headheight}{15pt}
\setlength{\headsep}{0.28cm}
\setlength{\footskip}{0.62cm}
\setlength{\parindent}{0pt}
\setlength{\parskip}{0.28em}
\setlength{\abovedisplayskip}{0.5em}
\setlength{\belowdisplayskip}{0.5em}
\setlength{\abovedisplayshortskip}{0.25em}
\setlength{\belowdisplayshortskip}{0.25em}
\setlist[enumerate]{leftmargin=2.2em,itemsep=0.3em,topsep=0.25em}
\setlist[itemize]{leftmargin=1.5em,itemsep=0.2em,topsep=0.2em}

\definecolor{repblue}{gray}{0}
\definecolor{repteal}{gray}{0}
\definecolor{repgray}{gray}{0}
\definecolor{repaccent}{gray}{0}
\definecolor{repline}{gray}{0}
\definecolor{replight}{gray}{1}

\hypersetup{
  hidelinks,
  pdftitle={TRGF Padilla Tarea 6},
  pdfauthor={Juan Ignacio Padilla Barrientos}
}

\pagestyle{fancy}
\fancyhf{}
\fancyhead[L]{\textcolor{repblue}{\small\itshape TRGF Padilla}}
\fancyhead[R]{\textcolor{repgray}{\small Seminario UCR I 2026}}
\fancyfoot[C]{\textcolor{repgray}{\small \thepage/\pageref*{LastPage}}}
\renewcommand{\headrulewidth}{0.3pt}
\renewcommand{\footrulewidth}{0pt}

\newcommand{\CC}{\mathbb{C}}
\newcommand{\RR}{\mathbb{R}}

\tcbset{
  enhanced,
  sharp corners,
  boxrule=0.45pt,
  arc=0pt,
  left=0.85em,
  right=0.85em,
  top=0.45em,
  bottom=0.4em,
  boxsep=0pt
}

\newcommand{\inlinehead}[2][repblue]{%
  \par\smallskip
  {\color{#1}\bfseries #2.}\hspace{0.45em}%
}

\newcommand{\answer}{\inlinehead{Respuesta}}

\newcommand{\topbanner}{%
  \begin{tcolorbox}[
    colback=replight,
    colframe=repline,
    borderline west={1.2pt}{0pt}{repblue},
    borderline east={1.2pt}{0pt}{repteal},
    left=1.0em,
    right=1.0em,
    top=0.7em,
    bottom=0.55em
  ]
    {\Large\bfseries\color{repblue} TRGF Padilla \textemdash\ Tarea 6}\hfill
    {\color{repaccent}\small\itshape Soluciones}\\[-0.1em]
    {\normalsize\color{repteal} Seminario UCR I 2026}
  \end{tcolorbox}
}

\newcounter{probcnt}
\newcount\workspacelines

\newcommand{\problem}[2]{%
  \stepcounter{probcnt}%
  \ifnum\value{probcnt}>1\newpage\fi
  \par\vspace{0.25em}
  {\large\bfseries\color{repblue} Problema #1}\par
  {\color{repline}\rule{\linewidth}{0.55pt}}\par
  \begin{tcolorbox}[
    breakable,
    colback=white,
    colframe=repline,
    borderline west={1.6pt}{0pt}{repteal},
    left=0.7em,
    right=0.7em,
    top=0.35em,
    bottom=0.3em,
    before skip=0.35em,
    after skip=0.45em
  ]
    #2
  \end{tcolorbox}
}

\newcommand{\workspace}[1]{%
  \answer
  \par\vspace{0.35em}
  \begingroup
  \workspacelines=0
  \loop
    \ifnum\workspacelines<#1
      \vspace{0.72cm}
      \advance\workspacelines by 1
  \repeat
  \endgroup
  \vspace{0.35em}
}

\begin{document}

\color{black}
\topbanner
\small

\problem{1}{%
Sea $(U,\pi)$ una representación de un grupo finito $G$. Suponga que $\chi_U(g) \in \RR$
para todo $g \in G$. Demuestre que existe un isomorfismo de representaciones
\[
  (U,\pi) \cong (U^*, \pi^*).
\]
}

\answer
Para la representación dual se tiene
\[
  \pi^*(g)(\lambda)=\lambda\circ \pi(g^{-1}),
  \qquad \lambda\in U^*.
\]
Por tanto, si $A_g$ es la matriz de $\pi(g)$ en alguna base de $U$, la matriz de
$\pi^*(g)$ en la base dual es $(A_g^{-1})^t$. Entonces
\[
  \chi_{U^*}(g)=\operatorname{tr}\bigl((A_g^{-1})^t\bigr)
  =\operatorname{tr}(A_g^{-1})=\chi_U(g^{-1}).
\]
Como $G$ es finito, $g$ tiene orden finito. Los valores propios de $A_g$ son raíces de
la unidad, de modo que los valores propios de $A_g^{-1}$ son los conjugados complejos de
los valores propios de $A_g$. Así,
\[
  \chi_U(g^{-1})=\overline{\chi_U(g)}.
\]
Por hipótesis $\chi_U(g)\in\RR$ para todo $g\in G$, luego
\[
  \chi_{U^*}(g)=\chi_U(g)
  \qquad (g\in G).
\]
Sobre $\CC$, dos representaciones de un grupo finito son isomorfas si y solo si tienen
el mismo carácter. Por consiguiente,
\[
  (U,\pi)\cong (U^*,\pi^*).
\]

\problem{2}{%
Sea $(U,\pi)$ una representación de un grupo $G$. Demuestre que $\chi_U(g)$ es un entero
algebraico para todo $g \in G$.

\smallskip
\textit{Puede utilizar que el conjunto de enteros algebraicos es cerrado bajo suma.}
}

\answer
Usaremos la hipótesis usual de este curso: los elementos considerados tienen orden finito
en la representación, en particular esto aplica si $G$ es finito. Fijemos $g\in G$ y sea
$m$ su orden. Entonces
\[
  \pi(g)^m=\pi(g^m)=\pi(e)=I.
\]
Así, el polinomio minimal de $\pi(g)$ divide a $x^m-1$. Como $x^m-1$ tiene raíces simples
en $\CC$, la matriz $\pi(g)$ es diagonalizable y todos sus valores propios son raíces
$m$-ésimas de la unidad. Si estos valores propios son $\lambda_1,\dots,\lambda_n$, entonces
\[
  \lambda_i^m=1
\]
para cada $i$. Por tanto, cada $\lambda_i$ es entero algebraico, pues es raíz del polinomio
mónico $x^m-1\in\mathbb Z[x]$.

Finalmente,
\[
  \chi_U(g)=\operatorname{tr}(\pi(g))=\lambda_1+\cdots+\lambda_n.
\]
Como los enteros algebraicos son cerrados bajo suma, $\chi_U(g)$ es un entero algebraico.

\smallskip
\textit{Nota de rigor.} Si $G$ fuera un grupo arbitrario sin hipótesis de finitud, el
enunciado no sería cierto en general. Por ejemplo, una representación de dimensión $1$ de
\(\mathbb Z\) puede enviar un generador a $1/2$, cuyo carácter no es entero algebraico.

\problem{3}{%
Sea ahora $G = S_4$.

\begin{enumerate}[label=(\alph*),leftmargin=2.4em]
\item (*) Demuestre que $S_4$ tiene $5$ clases conjugadas. Estas son:
\[
  \{e\},
\]
\[
  \{(12), (13), (14), (23), (24), (34)\},
\]
\[
  \{(12)(34), (13)(24), (14)(23)\},
\]
\[
  \{(123), (132), (134), (143), (124), (142), (234), (243)\},
\]
\[
  \{(1234), (1243), (1324), (1342), (1423), (1432)\}.
\]

\item (*) Demuestre que el subgrupo de $S_4$ dado por
\[
  H := \{e, (12)(34), (13)(24), (14)(23)\}
\]
es un subgrupo normal. Demuestre además que
\[
  S_4/H \cong S_3.
\]

\item (*) De la parte anterior, existe un morfismo de grupos
\[
  \varphi : S_4 \to S_3.
\]
Sean $V_{\mathrm{triv}}$ y $V_{\mathrm{sgn}}$ las dos representaciones irreducibles de dimensión
$1$ de $S_3$. Demuestre que, a través del homomorfismo $\varphi$, estas se pueden entender
como representaciones irreducibles de $S_4$. Verifique que los caracteres de estas dos
representaciones están dados por
\[
  \begin{array}{c|ccccc}
    & e & [(12)] & [(12)(34)] & [(123)] & [(1234)] \\
    \chi_{\mathrm{triv}} & 1 & 1 & 1 & 1 & 1 \\
    \chi_{\mathrm{sgn}} & 1 & -1 & 1 & 1 & -1
  \end{array}
\]
Pruebe que son ortogonales entre sí y que cada uno tiene norma $1$.

\item Sea $V_2$ la representación irreducible de dimensión $2$ de $S_3$. De nuevo, piense en
$V_2$ como una representación de $S_4$ por medio de $\varphi$. Demuestre que el carácter de
esta representación sería:
\[
  \begin{array}{c|ccccc}
    & e & [(12)] & [(12)(34)] & [(123)] & [(1234)] \\
    \chi_2 & 2 & 0 & 2 & -1 & 0
  \end{array}
\]
Verifique que la norma de $\chi_2$ es $1$ y, por lo tanto, $V_2$ es irreducible.

\item Sea $V$ la representación de $S_4$ de dimensión $4$ dada por permutar los elementos de
la base canónica $\{e_1, e_2, e_3, e_4\}$. Por ejemplo, en esta representación, las matrices
que representan a $(12)$, $(123)$ y $(13)(24)$ están dadas por
\[
  (12) \mapsto
  \begin{pmatrix}
    0 & 1 & 0 & 0 \\
    1 & 0 & 0 & 0 \\
    0 & 0 & 1 & 0 \\
    0 & 0 & 0 & 1
  \end{pmatrix},
  \qquad
  (123) \mapsto
  \begin{pmatrix}
    0 & 0 & 1 & 0 \\
    1 & 0 & 0 & 0 \\
    0 & 1 & 0 & 0 \\
    0 & 0 & 0 & 1
  \end{pmatrix},
\]
\[
  (13)(24) \mapsto
  \begin{pmatrix}
    0 & 0 & 1 & 0 \\
    0 & 0 & 0 & 1 \\
    1 & 0 & 0 & 0 \\
    0 & 1 & 0 & 0
  \end{pmatrix}.
\]
Demuestre que $V$ admite una subrepresentación de dimensión $1$ isomorfa a la representación
trivial. Sea $V_3$ el complemento ortogonal, de manera que
\[
  V \cong V_{\mathrm{triv}} \oplus V_3.
\]
Usando que $\chi_V = \chi_{\mathrm{triv}} + \chi_3$, demuestre que el carácter de $V_3$ está dado por
\[
  \begin{array}{c|ccccc}
    & e & [(12)] & [(12)(34)] & [(123)] & [(1234)] \\
    \chi_3 & 3 & 1 & -1 & 0 & -1
  \end{array}
\]
Demuestre que $\chi_3$ tiene norma $1$ y por lo tanto $V_3$ es irreducible. Verifique además
que $\chi_3$ es ortogonal a $\chi_{\mathrm{triv}}$, $\chi_{\mathrm{sgn}}$ y $\chi_2$.

\item (*) Escoja una base ortonormal para $V_3$ y determine las matrices que representan a
$(12)$, $(123)$ y $(1234)$ en esta base. Use esto para verificar de forma alternativa que el
carácter $\chi_3$ que calculamos arriba es correcto.

\item Sea $V_{3,\mathrm{sgn}}$ la representación de dimensión $3$ dada por
\[
  V_{3,\mathrm{sgn}} = V_3 \otimes V_{\mathrm{sgn}}.
\]
Utilice el ejercicio $2$ para verificar que el carácter de $V_{3,\mathrm{sgn}}$ está dado por
\[
  \begin{array}{c|ccccc}
    & e & [(12)] & [(12)(34)] & [(123)] & [(1234)] \\
    \chi_{3,\mathrm{sgn}} & 3 & -1 & -1 & 0 & 1
  \end{array}
\]
Verifique que $\chi_{3,\mathrm{sgn}}$ tiene norma $1$ y por tanto $V_{3,\mathrm{sgn}}$ es una
representación irreducible de $S_4$. Verifique que es ortogonal a $\chi_{\mathrm{triv}}$,
$\chi_{\mathrm{sgn}}$, $\chi_2$ y $\chi_3$.

\item Explique por qué $V_{\mathrm{triv}}$, $V_{\mathrm{sgn}}$, $V_2$, $V_3$ y
$V_{3,\mathrm{sgn}}$ clasifican todas las representaciones irreducibles de $S_4$ salvo
isomorfismo.

\item (*) Sea $W$ otra representación de $S_4$ tal que su carácter está dado por
\[
  \begin{array}{c|ccccc}
    & e & [(12)] & [(12)(34)] & [(123)] & [(1234)] \\
    \chi_W & 8 & 0 & 0 & -1 & 0
  \end{array}
\]
Determine cómo se descompone $W$ en suma directa de representaciones irreducibles.

\item (*) Determine cómo se descomponen los posibles tensores de irreps de $S_4$ como suma
de otras irreps.
\end{enumerate}
}

\answer
Trabajamos siempre en el orden de clases
\[
  e,\quad [(12)],\quad [(12)(34)],\quad [(123)],\quad [(1234)],
\]
cuyos tamaños son
\[
  1,\quad 6,\quad 3,\quad 8,\quad 6.
\]
Por tanto, si $\alpha$ y $\beta$ son caracteres de $S_4$, entonces
\[
  \langle \alpha,\beta\rangle
  =\frac{1}{24}\Bigl(
    \alpha_1\overline{\beta_1}
    +6\alpha_2\overline{\beta_2}
    +3\alpha_3\overline{\beta_3}
    +8\alpha_4\overline{\beta_4}
    +6\alpha_5\overline{\beta_5}
  \Bigr).
\]

\begin{enumerate}[label=(\alph*),leftmargin=2.4em]
\item En $S_n$, dos permutaciones son conjugadas si y solo si tienen el mismo tipo de
ciclo. Para $S_4$, las particiones de $4$ son
\[
  1+1+1+1,\qquad 2+1+1,\qquad 2+2,\qquad 3+1,\qquad 4.
\]
Estas dan exactamente cinco clases conjugadas. Sus elementos son, respectivamente:
\[
  \{e\},
\]
\[
  \{(12), (13), (14), (23), (24), (34)\},
\]
\[
  \{(12)(34), (13)(24), (14)(23)\},
\]
\[
  \{(123), (132), (134), (143), (124), (142), (234), (243)\},
\]
\[
  \{(1234), (1243), (1324), (1342), (1423), (1432)\}.
\]
Los tamaños se obtienen contando: hay $\binom42=6$ transposiciones, tres formas de partir
cuatro elementos en dos pares, $\binom43\cdot 2=8$ ciclos de longitud $3$, y $(4-1)!=6$
ciclos de longitud $4$.

\item El conjunto
\[
  H=\{e,(12)(34),(13)(24),(14)(23)\}
\]
es cerrado bajo producto: el producto de dos dobles transposiciones distintas es la tercera.
Así $H$ es un subgrupo de $S_4$, isomorfo al grupo de Klein. Además, $H$ es la unión de la
clase de $e$ con la clase de las dobles transposiciones. Como la conjugación preserva el
tipo de ciclo, para todo $\sigma\in S_4$ se cumple $\sigma H\sigma^{-1}=H$. Luego
$H\trianglelefteq S_4$.

Para identificar el cociente, consideremos las tres particiones de $\{1,2,3,4\}$ en dos
pares:
\[
  P_1=\bigl\{\{1,2\},\{3,4\}\bigr\},\quad
  P_2=\bigl\{\{1,3\},\{2,4\}\bigr\},\quad
  P_3=\bigl\{\{1,4\},\{2,3\}\bigr\}.
\]
Cada $\sigma\in S_4$ permuta el conjunto $\{P_1,P_2,P_3\}$, y por tanto define un
homomorfismo
\[
  \Phi:S_4\longrightarrow S(\{P_1,P_2,P_3\})\cong S_3.
\]
El núcleo de esta acción está formado por las permutaciones que fijan las tres particiones;
estas son precisamente
\[
  e,
  \qquad (12)(34),
  \qquad (13)(24),
  \qquad (14)(23).
\]
Así $\ker\Phi=H$. Entonces
\[
  |\operatorname{im}\Phi|=\frac{|S_4|}{|H|}=\frac{24}{4}=6=|S_3|,
\]
de modo que $\Phi$ es sobreyectivo. Por el primer teorema de isomorfismo,
\[
  S_4/H\cong S_3.
\]

\item Sea $\varphi:S_4\twoheadrightarrow S_3$ el epimorfismo anterior. Si
$(W,\rho)$ es una representación irreducible de $S_3$, entonces
$(W,\rho\circ\varphi)$ es irreducible como representación de $S_4$: cualquier subespacio
invariante bajo $\rho\circ\varphi(S_4)$ es invariante bajo $\rho(S_3)$ porque $\varphi$ es
sobreyectiva.

La imagen de las clases de $S_4$ bajo $\varphi$ tiene el siguiente tipo en $S_3$:
\[
  e\mapsto e,
  \qquad [(12)]\mapsto \text{transposiciones},
  \qquad [(12)(34)]\mapsto e,
\]
\[
  [(123)]\mapsto \text{$3$-ciclos},
  \qquad [(1234)]\mapsto \text{transposiciones}.
\]
Por tanto, al traer las representaciones trivial y signo de $S_3$ a $S_4$, obtenemos
\[
  \chi_{\mathrm{triv}}=(1,1,1,1,1),
  \qquad
  \chi_{\mathrm{sgn}}=(1,-1,1,1,-1).
\]
Sus normas son
\[
  \langle \chi_{\mathrm{triv}},\chi_{\mathrm{triv}}\rangle
  =\frac{1+6+3+8+6}{24}=1,
\]
\[
  \langle \chi_{\mathrm{sgn}},\chi_{\mathrm{sgn}}\rangle
  =\frac{1+6+3+8+6}{24}=1.
\]
Además,
\[
  \langle \chi_{\mathrm{triv}},\chi_{\mathrm{sgn}}\rangle
  =\frac{1-6+3+8-6}{24}=0.
\]

\item La representación irreducible de dimensión $2$ de $S_3$ tiene carácter
\[
  (2,0,-1)
\]
en las clases $e$, transposiciones y $3$-ciclos de $S_3$. Usando la descripción de las
imágenes de clases del inciso anterior, al verla como representación de $S_4$ obtenemos
\[
  \chi_2=(2,0,2,-1,0).
\]
Su norma es
\[
  \langle \chi_2,\chi_2\rangle
  =\frac{1\cdot 2^2+6\cdot 0^2+3\cdot 2^2+8\cdot (-1)^2+6\cdot 0^2}{24}
  =\frac{4+12+8}{24}=1.
\]
Por el criterio de irreducibilidad por caracteres, $V_2$ es irreducible como representación
de $S_4$.

\item En la representación de permutación $V=\CC^4$, el subespacio
\[
  L=\CC(e_1+e_2+e_3+e_4)
\]
es estable por $S_4$, y la acción sobre $L$ es trivial. Con el producto hermitiano usual,
su complemento ortogonal es
\[
  V_3=L^\perp
  =\{(z_1,z_2,z_3,z_4)\in\CC^4:z_1+z_2+z_3+z_4=0\}.
\]
Este subespacio también es estable por $S_4$, porque permutar coordenadas no cambia la suma.
Luego
\[
  V\cong V_{\mathrm{triv}}\oplus V_3.
\]

El carácter de $V$ cuenta puntos fijos de la permutación correspondiente. Así,
\[
  \chi_V=(4,2,0,1,0).
\]
Como $\chi_V=\chi_{\mathrm{triv}}+\chi_3$, se sigue que
\[
  \chi_3=(4,2,0,1,0)-(1,1,1,1,1)=(3,1,-1,0,-1).
\]
Entonces
\[
  \langle \chi_3,\chi_3\rangle
  =\frac{9+6+3+0+6}{24}=1.
\]
Por el criterio de irreducibilidad, $V_3$ es irreducible. Además,
\[
  \langle \chi_3,\chi_{\mathrm{triv}}\rangle
  =\frac{3+6-3-6}{24}=0,
\]
\[
  \langle \chi_3,\chi_{\mathrm{sgn}}\rangle
  =\frac{3-6-3+6}{24}=0,
\]
y
\[
  \langle \chi_3,\chi_2\rangle
  =\frac{6-6}{24}=0.
\]

\item Tomemos la base ortonormal de $V_3$ dada por
\[
  b_1=\frac{e_1-e_2}{\sqrt2},
  \qquad
  b_2=\frac{e_1+e_2-2e_3}{\sqrt6},
  \qquad
  b_3=\frac{e_1+e_2+e_3-3e_4}{\sqrt{12}}.
\]
En esta base, las matrices de $(12)$, $(123)$ y $(1234)$ son
\[
  [(12)]_{\mathcal B}
  =
  \begin{pmatrix}
    -1 & 0 & 0 \\
    0 & 1 & 0 \\
    0 & 0 & 1
  \end{pmatrix},
\]
\[
  [(123)]_{\mathcal B}
  =
  \begin{pmatrix}
    -\frac12 & -\frac{\sqrt3}{2} & 0 \\
    \frac{\sqrt3}{2} & -\frac12 & 0 \\
    0 & 0 & 1
  \end{pmatrix},
\]
y
\[
  [(1234)]_{\mathcal B}
  =
  \begin{pmatrix}
    -\frac12 & -\frac{\sqrt3}{6} & -\frac{\sqrt6}{3} \\
    \frac{\sqrt3}{2} & -\frac16 & -\frac{\sqrt2}{3} \\
    0 & \frac{2\sqrt2}{3} & -\frac13
  \end{pmatrix}.
\]
Sus trazas son, respectivamente,
\[
  1,
  \qquad 0,
  \qquad -1.
\]
Además, la identidad tiene traza $3$ y una doble transposición tiene traza $-1$ en $V_3$
porque, por el inciso anterior, $\chi_3((12)(34))=\chi_V((12)(34))-1=0-1=-1$.
Esto recupera
\[
  \chi_3=(3,1,-1,0,-1).
\]

\item Como
\[
  V_{3,\mathrm{sgn}}=V_3\otimes V_{\mathrm{sgn}},
\]
su carácter es el producto punto a punto de los caracteres:
\[
  \chi_{3,\mathrm{sgn}}(g)=\chi_3(g)\chi_{\mathrm{sgn}}(g).
\]
Por tanto,
\[
  \chi_{3,\mathrm{sgn}}
  =(3,1,-1,0,-1)(1,-1,1,1,-1)
  =(3,-1,-1,0,1).
\]
Su norma es
\[
  \langle \chi_{3,\mathrm{sgn}},\chi_{3,\mathrm{sgn}}\rangle
  =\frac{9+6+3+0+6}{24}=1,
\]
luego $V_{3,\mathrm{sgn}}$ es irreducible. También se verifica
\[
  \langle \chi_{3,\mathrm{sgn}},\chi_{\mathrm{triv}}\rangle
  =\frac{3-6-3+6}{24}=0,
\]
\[
  \langle \chi_{3,\mathrm{sgn}},\chi_{\mathrm{sgn}}\rangle
  =\frac{3+6-3-6}{24}=0,
\]
\[
  \langle \chi_{3,\mathrm{sgn}},\chi_2\rangle
  =\frac{6-6}{24}=0,
\]
y
\[
  \langle \chi_{3,\mathrm{sgn}},\chi_3\rangle
  =\frac{9-6+3-6}{24}=0.
\]

\item Hemos construido cinco representaciones irreducibles no isomorfas de $S_4$:
\[
  V_{\mathrm{triv}},\quad V_{\mathrm{sgn}},\quad V_2,\quad V_3,
  \quad V_{3,\mathrm{sgn}}.
\]
Son no isomorfas porque sus caracteres son distintos y ortogonales. Además, $S_4$ tiene
exactamente cinco clases conjugadas, y el número de representaciones irreducibles complejas
de un grupo finito es igual al número de clases conjugadas. Por tanto estas cinco clasifican
todas las irreducibles de $S_4$ salvo isomorfismo.

También se puede ver por la suma de cuadrados de dimensiones:
\[
  1^2+1^2+2^2+3^2+3^2=24=|S_4|.
\]

\item Si $W$ tiene carácter
\[
  \chi_W=(8,0,0,-1,0),
\]
la multiplicidad de una irrep de carácter $\chi$ en $W$ es
\[
  \langle \chi_W,\chi\rangle.
\]
Calculamos:
\[
  \langle \chi_W,\chi_{\mathrm{triv}}\rangle
  =\frac{8-8}{24}=0,
  \qquad
  \langle \chi_W,\chi_{\mathrm{sgn}}\rangle
  =\frac{8-8}{24}=0,
\]
\[
  \langle \chi_W,\chi_2\rangle
  =\frac{16+8}{24}=1,
\]
\[
  \langle \chi_W,\chi_3\rangle
  =\frac{24}{24}=1,
  \qquad
  \langle \chi_W,\chi_{3,\mathrm{sgn}}\rangle
  =\frac{24}{24}=1.
\]
Así,
\[
  W\cong V_2\oplus V_3\oplus V_{3,\mathrm{sgn}}.
\]

\item Para descomponer productos tensoriales usamos que
\[
  \chi_{A\otimes B}(g)=\chi_A(g)\chi_B(g)
\]
y luego tomamos productos internos con los cinco caracteres irreducibles. Como el producto
tensorial es conmutativo salvo isomorfismo, basta listar los productos siguientes:
\[
  V_{\mathrm{triv}}\otimes X\cong X
  \qquad\text{para toda irrep }X,
\]
\[
  V_{\mathrm{sgn}}\otimes V_{\mathrm{sgn}}\cong V_{\mathrm{triv}},
  \qquad
  V_{\mathrm{sgn}}\otimes V_2\cong V_2,
  \qquad
  V_{\mathrm{sgn}}\otimes V_3\cong V_{3,\mathrm{sgn}},
\]
\[
  V_{\mathrm{sgn}}\otimes V_{3,\mathrm{sgn}}\cong V_3.
\]
Los productos restantes son
\[
  V_2\otimes V_2
  \cong V_{\mathrm{triv}}\oplus V_{\mathrm{sgn}}\oplus V_2,
\]
\[
  V_2\otimes V_3
  \cong V_3\oplus V_{3,\mathrm{sgn}},
  \qquad
  V_2\otimes V_{3,\mathrm{sgn}}
  \cong V_3\oplus V_{3,\mathrm{sgn}},
\]
\[
  V_3\otimes V_3
  \cong V_{\mathrm{triv}}\oplus V_2\oplus V_3\oplus V_{3,\mathrm{sgn}},
\]
\[
  V_3\otimes V_{3,\mathrm{sgn}}
  \cong V_{\mathrm{sgn}}\oplus V_2\oplus V_3\oplus V_{3,\mathrm{sgn}},
\]
\[
  V_{3,\mathrm{sgn}}\otimes V_{3,\mathrm{sgn}}
  \cong V_{\mathrm{triv}}\oplus V_2\oplus V_3\oplus V_{3,\mathrm{sgn}}.
\]
Esto determina todos los tensores posibles de irreps de $S_4$.
\end{enumerate}

\end{document}